\documentclass[10pt,a4paper]{amsart}
\usepackage[margin=19mm]{geometry}
\usepackage{amsmath,amssymb,mathtools}
\usepackage[scr=rsfs,cal=euler]{mathalfa}
\usepackage{xfrac}
\usepackage[unicode,hidelinks,bookmarks=false]{hyperref}
\newcommand{\ie}{i.e.\ }
\newcommand{\cf}{cf.\ }
\newcommand{\resp}{respectively\ }
\newcommand{\Subsection}{\subsection}
\providecommand{\nobreakdash}{\nobreak\hbox{-}}
\setlength{\emergencystretch}{2em}
\multlinegap0pt
\usepackage{amssymb}
\usepackage{breqn}
\usepackage{nicefrac}
\usepackage[all]{xy}
\usepackage{enumerate}
\usepackage{mathtools}
\usepackage{float}
\usepackage{caption}
\usepackage{tikz}
\newtheorem{proposition}{Proposition}
\newtheorem{lemma}{Lemma}
\newcommand{\RR}{{\mathbb R}}
\title{Symplectic variations of convex bodies and the mean width}
\author{Jonghyeon Ahn, Ely Kerman}
\date{Source: arXiv:2311.02870v3 (CC BY 4.0)}
\begin{document}
\maketitle

\noindent\emph{Source and licence.} Symplectic variations of convex bodies and the mean width. Version arXiv:2311.02870v3 (CC BY 4.0), distributed by arXiv under the Creative Commons Attribution 4.0 International licence, http://creativecommons.org/licenses/by/4.0/, as recorded in the arXiv licence metadata for this exact version and shown on the abstract page https://arxiv.org/abs/2311.02870v3. Full licence notice: \texttt{licenses/arxiv-2311.02870-CC-BY-4.0.txt}.\par\smallskip

\noindent\textit{Editorial scope.} Counted unit: the article's lemma nail (source0.tex lines 1286--1296). The article's complete original proof of it is delivered with it (source0.tex lines 1298--1322, one \texttt{proof} environment of 25 lines). No mathematics is written, repaired or re-derived: the statement and the proof are the article's own lines, in the article's own notation, and no retained line is substituted or re-flowed.  Retained uncounted as the source-local context the proof needs: lines 1043--1049; lines 1202--1204; lines 1221--1223.  Nothing was omitted from any retained span, so the package carries no omission record.  No retained line contains a citation, so the wrapper carries no bibliography and no bibliography entry.  No retained line contains a figure environment, an image-inclusion command or drawing code, so no image is delivered and the package declares no asset dependency.  Editorial formatting only, in addition: an AMS \texttt{amsart} preamble that restates the article's own declarations and macros used by the retained lines, namely the environments \texttt{proposition}, \texttt{lemma} on separate counters, together with the standard packages those lines need.  The retained environments print in this document as Proposition 1, Lemma 1 in order of appearance, whereas the article prints the same items with the numbering of its own full document.  The complete native source of the article as distributed in the arXiv e-print for this exact version is bundled unchanged as \texttt{sources/149582/source0.tex}.
\begin{proposition}\label{form}
Let $K =\{\rho=f(r)\}$ be a strictly convex toric domain in $\RR^{2n}$ with smooth boundary. For any smooth function $H \colon \RR^{2n} \to \RR$ the function
\begin{align*}
\left\langle u, \left.\frac{d}{dt}\right|_{t=0} \nu^{-1}_{K_t}(u)\right\rangle
\end{align*}
is $\theta$-simple where $K_t =\phi^t_H(K)$.
\end{proposition}

\begin{lemma}\label{Dv}
    The function $$u=(\theta,r) \mapsto \left\langle u, Dw_K^{-1}(u)\left[  v_0(w_K^{-1}(u))\right] \right\rangle$$ only depends on $r$.
\end{lemma}

\begin{lemma}\label{Dvdot}
    The function $$u=(\theta,r) \mapsto \left\langle u, Dw_K^{-1}(u)\left[  \dot{v}_0(w_K^{-1}(u))\right] \right\rangle$$ is $\theta$-simple for any smooth $H \colon \RR^{2n}\to \RR$.
\end{lemma}

\begin{lemma}\label{nail}
The function 
\begin{align*}
     df_0 (w_{K}^{-1}(u)) \left[Dw_K^{-1}(u)\left[ v_0(w^{-1}_{K}(u))\right]\right]
\end{align*}
only depends on $r$, and the function
\begin{align*}
     df_0 (w_{K}^{-1}(u)) \left[Dw_K^{-1}(u)\left[ \dot{v}_0(w^{-1}_{K}(u))\right]\right]
\end{align*}
is $\theta$-simple.
\end{lemma}

\begin{proof}
    In the notation of the proof of Lemma \ref{Dv}, we have
    \begin{align*}Dw_K^{-1}(u)\left[  v_0(w_K^{-1}(u))\right] =& \sum_{j=1}^n \left( r_j\sum_{k=1}^{n-1}C_k(r)\frac{\partial a_j}{\partial r_k}(r)+ a_j(r)C_j(r) \right)\cos \theta_j \frac{\partial}{\partial x_j}\\ &+\nonumber \sum_{j=1}^n \left( r_j\sum_{k=1}^{n-1}C_k(r)\frac{\partial a_j}{\partial r_k}(r)+ a_j(r)C_j(r) \right)\sin \theta_j \frac{\partial}{\partial y_j}.\end{align*}
Since
\begin{align*}
    df_0(u) = \sum_{j=1}^{n-1} \frac{\partial f_0}{\partial r_j}(r) dr_j =\sum_{j=1}^{n-1} \frac{\partial f_0}{\partial r_j}(r) (\cos \theta_j\, dx_j + \sin \theta_j \,dy_j),
\end{align*}
we have
\begin{align*}
    df_0(w_K^{-1}(u)) =\sum_{j=1}^{n-1} \frac{\partial f_0}{\partial r_j}(F(r)) (\cos \theta_j\, dx_j + \sin \theta_j \,dy_j).
\end{align*}
Hence,  
\begin{align*}
     df_0 (w_{K}^{-1}(u)) \left[Dw_K^{-1}(u)\left[ v_0(w^{-1}_{K_t}(u))\right]\right] = \sum_{j=1}^n \frac{\partial f_0}{\partial r_j}(F(r))\left( r_j\sum_{k=1}^{n-1}C_k(r)\frac{\partial a_j}{\partial r_k}(r)+ a_j(r)C_j(r) \right)
\end{align*}
and the first assertion of Lemma has been verified. To verify the second, we use the notation of the proof of Lemma \ref{Dvdot} to write
\begin{align*}
 &Dw_K^{-1}(u)\left[  \dot{v}_0(w_K^{-1}(u))\right] \\=& \sum_{j=1}^n \left( a_j(r)A_r(\theta,r)\sin \theta_j  + \left(r_j\sum_{k=1}^{n-1}B_k(\theta,r) \frac{\partial a_j}{\partial r_k}(r) +a_j(r)B_j(\theta,r)\right)\cos \theta_j  \right)\frac{\partial}{\partial x_j}\\&+ \sum_{j=1}^n \left( -a_j(r)A_r(\theta,r)\cos \theta_j  + \left(r_j\sum_{k=1}^{n-1}B_k(\theta,r) \frac{\partial a_j}{\partial r_k}(r) +a_j(r)B_j(\theta,r)\right)\sin \theta_j\right)\frac{\partial}{\partial y_j}.
\end{align*}
It follows that 
\begin{align*}
     df_0 (w_{K}^{-1}(u)) \left[Dw_K^{-1}(u)\left[ \dot{v}_0(w^{-1}_{K_t}(u))\right]\right] = \sum_{j=1}^n  \frac{\partial f_0}{\partial r_j}(F(r))\left(r_j\sum_{k=1}^{n-1}B_k(\theta,r) \frac{\partial a_j}{\partial r_k}(r) +a_j(r)B_j(\theta,r)\right).
\end{align*}
Since the $B_j(\theta,r)$ are $\theta$-simple, the proof of Proposition \ref{form} is complete.
\end{proof}

\end{document}
